% META: {"id": "TCOMPL-CORR-EX-027", "chapitre": "TCOMPL-CORRELATION-CAUSALITE", "type_objet": "exercice", "capacites": ["TCOMPL-CORR-C2"], "capacites_codes": ["C2"], "methodes": ["M4"], "parcours": 2, "competences": ["communiquer","calculer"], "duree_min": 10, "mode_creation": "ex_nihilo", "sources_inspiration": [], "fichier_tex": "chapitres/TCOMPL-CORRELATION-CAUSALITE/exercices/TCOMPL-CORR-EX-027.tex", "corrige_tex": "chapitres/TCOMPL-CORRELATION-CAUSALITE/corriges/TCOMPL-CORR-CO-027.tex", "status": "approved"}
% BEGIN-VERIFY
% from sympy import *
% x = symbols('x')
% # f(x)=x^4-2x^2, f'=4x^3-4x, f''=12x^2-4
% f = x**4 - 2*x**2
% fp = diff(f, x)
% fpp = diff(f, x, 2)
% assert expand(fp) == 4*x**3 - 4*x
% assert expand(fpp) == 12*x**2 - 4
% assert solve(fpp, x) == [-sqrt(3)/3, sqrt(3)/3]
% assert f.subs(x, 0) == 0
% assert f.subs(x, sqrt(3)/3) == Rational(-5, 9)
% END-VERIFY
\begin{exercice}{TCOMPL-CORR-EX-027}{1}{10}
Soit $f(x) = x^4 - 2x^2$.
\begin{enumerate}
  \item Calculer $f'(x)$ et $f''(x)$.
  \item Determiner les points d'inflexion de la courbe.
  \item Esquisser l'allure de la courbe.
\end{enumerate}

\end{exercice}
